\documentclass[11pt]{article}
\usepackage[a4paper,margin=2.6cm]{geometry}
\usepackage{amsmath,amssymb,amsthm}
\title{The spectral square root of $Q_L$\\
\large high-pass on the evaluators, crush on the desert}
\author{D.\ Joubert}
\date{4 September 2026}
\begin{document}
\maketitle
\begin{abstract}
\noindent
$Q^{1/2}=\sum\sqrt{\lambda_j}\,v_jv_j^{\!\top}$ is the unique PSD
factor. Its top eigenspace is the span of the in-band evaluators
$\hat\eta(\gamma_k)$ (1--1 pairing for $K\le6$). Its bottom eigenspace
is the radical: $\sqrt{\lambda_{\min}}\sim10^{-24}$ at $\mu=11$. As an
operator on $V_N$, $Q^{1/2}$ is a high-pass tuned to the zeros, not a
list of Euler lags. Polar of the evaluator matrix $C$ recovers
$Q^{1/2}$ only after the off-band tail is in; with four in-band zeros
on $V_{13}$ the gap is $47\%$. No new closed form.
\end{abstract}

\section{The operator}
On $V_N$,
\[
Q^{1/2}f=\sum_{j=1}^{N}\sqrt{\lambda_j}\,\langle f,v_j\rangle\,v_j.
\]
If $C$ is the matrix with columns $\hat\eta(\gamma_k)$ then the
explicit formula says $Q\approx CC^{\!\top}$. The left polar factor of
$C$ is $(CC^{\!\top})^{1/2}$. Hence
\[
Q^{1/2}\;\approx\;(CC^{\!\top})^{1/2}
\]
as soon as the tail of $\gamma$ beyond $\omega_{\max}$ is small.
At $\mu=11$, $\|Q-Q_K\|_F/\|Q\|_F$ is $8.9\%$ at $K=40$
($\gamma\approx\omega_{\max}$) and $5\%$ at $K=60$ (notebook P1). The
same tail sits in the polar: $Q^{1/2}$ and $(CC^{\!\top})^{1/2}$
differ by that remainder, not by a new operator.

\section{Two regimes}
\paragraph{High.}
Greedy pairing of the largest $v_j$ with the in-band $\hat\eta(\gamma_k)$
is 1--1 for $K\le6$ (mean cosine $0.81$ at $N=13$, $K=4$). It breaks
at the last in-band zero ($K=7$, min cosine $0.27$). Those
$\sqrt{\lambda_j}$ are $O(1)$ (on $V_{13}$: $1.74,1.59,1.56,1.42$).
$Q^{1/2}$ acts on that subspace as a mild gain.

\paragraph{Low.}
The remaining eigenvalues follow the Slepian staircase of the desert
and then the razor. On $V_{13}$ already $\sqrt{\lambda}\approx
0.27,0.024,0.001,\ldots$. On $V_{47}$, $\sqrt{\lambda_{\min}}\approx
6\times10^{-24}$. $Q^{1/2}$ crushes the ground state. That is the
same hyper-nullity as MUSIC: the radical is the approximate kernel of
every in-band evaluator, hence of $C^{\!\top}$, hence of $Q^{1/2}$.

\section{What this is not}
$Q^{1/2}$ is not an arithmetic $T_L$. Its rows in the cosine basis
are the $v_j$ weighted by $\sqrt{\lambda_j}$, and the $v_j$ of the
top are linear combinations of evaluators at the $\gamma_k$, not at
the $\log p$. Passing from $Q^{1/2}$ to a factor whose rows are lags
$\tau_{\log p}$ is the SOS problem already recorded as open
(\emph{sos-arithmetic}). The spectral square root is the unique PSD
answer to ``write $Q$ as a norm''; it answers that question by
pointing back at the zeros.

\section*{Status}
Analyzed. $Q^{1/2}$ is the polar of the evaluator synthesis, up to
the off-band tail. High-pass on $\mathrm{span}\{\hat\eta(\gamma_k)\}$,
crush on the desert. Not a closed arithmetic factor. Not RH.
\end{document}
